\chapter{Geometry of Laws and Local Consistency}\label{chap:2}

\section*{Overview of the Chapter}

The question of this chapter is \textbf{when states that satisfy the Laws on each part can
be glued into a state of the whole}. Even if the Laws hold on each part, the parts need not
represent the same state once they are connected. Suppose, for example, that three services
record the same value against different references. Even if a correction between references
can be specified for each connection, no reference common to all three can be chosen unless
the corrections add up to zero around the cycle through the three. This discrepancy is
invisible from the equations inside the individual services alone.

This chapter treats the problem in two stages. First, we collect the equations of the Laws
into an ideal and construct the space of states that satisfy the equations. Second, we
compare the differences of the states chosen on the parts over the overlaps, and remove the
differences that can be resolved by local corrections. The equivalence class that remains
is the obstruction class that prevents gluing.

\begin{xltabular}{\linewidth}{@{}LLL@{}}
\toprule
Question & Construction & Outcome \\
\midrule
\endhead
Where do the states that satisfy the Laws live? & Construct coordinate rings, ideals of equations, and evaluations & Equivalence of the Laws holding with factorization through the lawful locus (\S\S\ref{sec:2.1}--\ref{sec:2.2}) \\
\addlinespace[3pt]
What is still missing after local success? & Construct differences on overlaps and the \v{C}ech complex & The obstruction class of specified local data, and gluing conditions after correction (\S\S\ref{sec:2.3}--\ref{sec:2.5}) \\
\addlinespace[3pt]
How do covers and coefficients affect obstructions? & Examine cycles, dimensions, differences on boundaries, and the structural and semantic phases & Conditions for vanishing, detection, and comparison (\S\S\ref{sec:2.6}--\ref{sec:2.7}) \\
\addlinespace[3pt]
How are concrete obstructions passed to the diagnosis of the next chapter? & Build integer-coefficient local data from Atoms, generator relations, and an actual cover & Obstruction classes computed from states and transitions (\S\ref{sec:2.8}) \\
\addlinespace[3pt]
Do the semantics of repair and the computation with equations agree? & Construct and compare two coefficients and their local states & Correspondence of the SAGA obstruction classes, and conditions for the existence of a global repair (\S\S\ref{sec:2.9}--\ref{sec:2.10}) \\
\bottomrule
\end{xltabular}

The central conclusions are that, under the necessary conditions on evaluation, the
validity of the Laws can be read as a geometric zero-locus condition, and that, when the
actual repair states form a sheaf, a global repair is obtained from the vanishing of the
obstruction class. The former requires a map connecting the symbolic equations with their
evaluation, and the latter requires the action of local corrections and gluing. We define
each of these maps and conditions before the theorems that use them.

\subsection*{Why we connect to the semantics of repair}

In Chapter~\ref{chap:1} we encoded the independently defined semantics of lenses and
protocols into objects of AAT, and recovered the operations and the semantics-preserving
morphisms from those objects and typed morphisms. We also showed that the validity of the
constructed Laws is equivalent to the validity of the laws of the original semantics. For
local consistency as well, we need to know how the computed obstructions correspond to the
repairs that are actually permitted. To this end, we specify which local changes count as
the same repair and which residuals of the equations are identified, and prove the
correspondence between the two.

This comparison is given by SAGA, the preceding paper
(\cite[\href{https://arxiv.org/abs/2608.21458}{\S\S4--5}]{SAGA}).
In this chapter we construct its coefficients, local states, and comparison maps, and prove
the correspondence of obstruction classes and the existence of repairs. This lets us carry
the zero and nonzero verdicts obtained on the equation side back to repairability in the
chosen semantics.

\section{Local algebras and ideals of equations}\label{sec:2.1}

To solve equations, we first fix the variables and the structural relations that they
satisfy from the outset. Using the ring attached to each context in Chapter~\ref{chap:1},
we incorporate, in order, the relations expressing the structure and the equations of the
Laws to be checked.

Throughout, we fix a reading of the architecture, a site $(\mathcal{C},J)$, and an equation
system $E$. We write $W$ for a context, $i\in K_E$ for a Law index, and
$K_E^{\mathrm{req}}$ for the set of required indices. The two families of coordinates
$\nu_{W,i,a}$ and $\varepsilon_{W,A,i,a}$ of Definition~\ref{def:1.12} represent the
symbolic equations and the residuals at an object $A$, respectively.

\begin{definition}[Presentation of local algebras]\label{def:2.1}

For the coordinate set $Z_W$ and the ideal of structural relations $I_W^{\mathrm{str}}$ of
Definition~\ref{def:1.25}, we choose a $k$-algebra isomorphism
\begin{equation*}
B_{\mathrm{raw}}(W)=k[x_z\mid z\in Z_W]/I_W^{\mathrm{str}},
\qquad
\theta_W:B_{\mathrm{raw}}(W)\xrightarrow{\sim}O_E(W)
\tag{2.1}\label{eq:2.1}
\end{equation*}
We require that $\theta_W$ commute with the restrictions of contexts and send the equation
coordinates and the residual coordinates in the presentation to the corresponding $\nu$ and
$\varepsilon$. We call this a presentation of the local algebras of $E$. The sheaf of rings
$\mathcal{O}=a_JO_E$ is obtained by transferring addition and multiplication to the
sheafification of Proposition~\ref{prop:1.23}. The operations can be computed on the local
presentation, and the ring axioms are preserved as local identities.

\end{definition}

\begin{proposition}[Algebraic presentation of local configurations]\label{prop:2.2}

For a commutative $k$-algebra $R$, let $\mathrm{Conf}_W(R)$ be the set of families of
coordinate values $(c_z)_{z\in Z_W}$ that satisfy all the structural relations. Then there
is a bijection
\[
\mathrm{Conf}_W(R)
\cong\mathrm{Hom}_{k\text{-}\mathrm{Alg}}(B_{\mathrm{raw}}(W),R)
\]
natural in $R$.

\end{proposition}

\begin{proof}

A family of coordinate values determines a unique homomorphism from the polynomial ring to
$R$. That all structural relations map to zero is equivalent to the kernel containing
$I_W^{\mathrm{str}}$, so the homomorphism factors through the quotient. Conversely,
evaluating a homomorphism from the quotient at each $x_z$ recovers the original family.
Composition with a ring homomorphism $R\to R'$ is exactly the transfer of coordinate
values, so naturality follows.

\end{proof}

The relations imposed here are the structural relations that describe one and the same
local configuration. The Laws that such a configuration must satisfy are expressed by the
following ideals.

\begin{definition}[Witness ideals and obstruction ideals]\label{def:2.3}

We define the witness ideal of a Law $i$ and the obstruction ideal of the required Laws by
\begin{equation*}
\begin{aligned}
I_i^E(W)&=(\nu_{W,i,a}\mid a\in\mathrm{At})\subseteq O_E(W),\\
I_{\mathrm{Ob}}^E(W)&=\sum_{i\in K_E^{\mathrm{req}}}I_i^E(W)
\end{aligned}
\tag{2.2}\label{eq:2.2}
\end{equation*}
An element of a generated ideal is a finite sum of the generators with ring coefficients.
Hence, even when the index set is infinite, each individual element is expressed by a
finite sum.

\end{definition}

\begin{proposition}[Restriction and sheafification of ideals]\label{prop:2.4}

A restriction $j:W'\to W$ satisfies $\mathrm{res}_j(I_i^E(W))\subseteq I_i^E(W')$. The same
holds for $I_{\mathrm{Ob}}^E$. The image of the sheafified morphism
$a_JI_i^E\to\mathcal{O}$ is a sheaf of ideals of $\mathcal{O}$. We write it
$\mathcal{I}_i^E$, and write $\mathcal{I}_{\mathrm{Ob}}^E$ for its sum over the required
indices.

\end{proposition}

\begin{proof}

The restriction of a generator is $\nu_{W',i,a}$, and the restriction of any element of the
generated ideal is again a finite sum of generators. Closure under addition and under
multiplication by ring elements can be checked locally, so it is preserved in the image
after sheafification. For the sum, the same argument applies once each section is expressed
locally as a finite sum.

\end{proof}

Enlarging the set of required indices enlarges the ideal and strengthens its zero-locus
condition. When optional Laws are imposed at the same time, their ideals are added to the
sum.

\section{Spaces satisfying the Laws and evaluation}\label{sec:2.2}

An ideal collects the equations, but by itself it does not express that the equations hold
at an object. In this section we read the rings as geometric spaces and provide the
mechanism by which the symbolic equations are evaluated into the residuals of objects.

\begin{definition}[Affine schemes and closed subschemes]\label{def:2.5}

For a commutative ring $B$, $\mathrm{Spec}\,B$ is the affine scheme whose underlying set is
the set of prime ideals, and whose basic open sets and their rings are given by
\begin{equation*}
D(f)=\{\mathfrak{p}\mid f\notin\mathfrak{p}\},
\qquad
\mathcal{O}_{\mathrm{Spec}\,B}(D(f))=B_f
\tag{2.3}\label{eq:2.3}
\end{equation*}
Here $B_f$ is the localization that inverts $f$. Sections over a general open set are
defined by gluing compatible sections over basic open sets. The local ring at each point is
$B_{\mathfrak{p}}$. A space equipped with a sheaf of rings that locally has this form is
called a scheme. These are the usual constructions of the prime spectrum and the structure
sheaf
(\cite[\href{https://stacks.math.columbia.edu/tag/01HR}{\S26.5, Tag~01HR}]{Stacks}).

We write $V_B(I)=\mathrm{Spec}(B/I)$ for the closed subscheme defined by an ideal
$I\subseteq B$. A ring homomorphism $e:B\to R$ factors through $B/I$ if and only if
$e(I)=0$. Hence $V_B(I)$ represents the points that satisfy the equations $I$ in an
arbitrary coefficient ring.

\end{definition}

\subsection*{Evaluating symbolic equations into residuals}

We impose conditions in two stages on a space representing local evaluation data. First, we
make the evaluation of the symbolic equations agree with the evaluation of the residuals of
objects. Second, we set the required residuals to zero. Below, $Y_W$, $X_W$, and
$X_E^{\mathrm{law}}\cap X_W$ denote, in order, the space before the conditions are imposed
and the spaces after each condition is imposed.

\begin{construction}[Geometric realization of the equations]\label{cons:2.6}

In each context, consider local evaluation data $p$ consisting of a reading $A_p$ of an
object and a homomorphism $e_p:O_E(W)\to R$ into a commutative $k$-algebra $R$. We assume
that these data can be transferred along changes of the coefficient ring. We treat the case
in which the following algebraic presentation is given.

\begin{enumerate}
\item The functor of local evaluation data is represented by an affine scheme
$Y_W=\mathrm{Spec}\,D_W$, and a universal ring homomorphism $u_W:O_E(W)\to D_W$ gives the
evaluation.
\item The evaluation of the residuals $p\mapsto e_p(\varepsilon_{W,A_p,i,a})$ is a regular
function, and an element $r_{W,i,a}\in D_W$ representing it has been constructed.
\item The comparisons of the presentations between contexts give isomorphisms on open
parts and preserve structure, coordinates, and residuals on the overlaps. On triple
overlaps the open parts being compared correspond, and the composites of the isomorphisms
agree there.
\item The assignment of charts after gluing, $W\mapsto X_W$ (defined in \eqref{eq:2.4}
below), sends morphisms of contexts to inclusions of open parts and preserves identities,
composites, and overlaps. Covers in $J$ are sent to open covers of charts.
\end{enumerate}

The first condition says that the evaluation data can be described by coordinates, and the
second that the residuals are algebraic functions of those coordinates. The third is the
condition for gluing the local presentations without contradiction. The fourth says that
restrictions, overlaps, and covers of contexts can be handled in the same way on the open
parts of the resulting space. We then impose the relations that equate the evaluation of
the symbolic equations with that of the residuals.
\begin{equation*}
\begin{aligned}
K_W&=(u_W(\nu_{W,i,a})-r_{W,i,a}\mid i,a),\\
B_W&=D_W/K_W,\qquad X_W=\mathrm{Spec}\,B_W,\\
\eta_W&:O_E(W)\longrightarrow B_W.
\end{aligned}
\tag{2.4}\label{eq:2.4}
\end{equation*}
Since the comparison isomorphisms preserve $K_W$, the open parts of the $X_W$ can also be
compared by isomorphisms. Let $X_E$ be the scheme obtained by gluing them. Concretely, we
identify corresponding points in the disjoint union of the $X_W$, and define the functions
on an open set as the families of functions that agree on each chart. The composition
condition on the overlaps makes the identification transitive, and each $X_W$ is embedded
as an open part. The local rings are also those of the charts, so the resulting space is a
scheme
(\cite[\href{https://stacks.math.columbia.edu/tag/01JA}{\S26.14, Tag~01JA}]{Stacks}).

For $s:T\to X_E$, put $T_W=T\times_{X_E}X_W$. We define its evaluation
$\mathrm{ev}_{s,W}:O_E(W)\to\Gamma(T_W,\mathcal{O}_T)$ as the composite of $\eta_W$ with
the pullback of regular functions. We write $A_s$ for the object read off from the
representing functor, and $\varepsilon_{W,i,a}(s)=\mathrm{ev}_{s,W}(\varepsilon_{W,A_s,i,a})$
for its residuals. For a general $T$, we read this formula on affine open sets and glue the
agreeing values.

\end{construction}

\begin{lemma}[Evaluation of generators and localization]\label{lem:2.7}

The evaluation obtained in Construction~\ref{cons:2.6} is compatible with restriction of
contexts and with base change of $T$, and satisfies
\begin{equation*}
\mathrm{ev}_{s,W}(\nu_{W,i,a})=\varepsilon_{W,i,a}(s)
\tag{2.5}\label{eq:2.5}
\end{equation*}
Moreover, putting $J_i(W)=(\eta_W(\nu_{W,i,a})\mid a)$, on a basic open set
$D(f)\subseteq X_W$ this ideal becomes
\begin{equation*}
J_i(W)B_{W,f}
=(\eta_W(\nu_{W,i,a})/1\mid a)
\tag{2.6}\label{eq:2.6}
\end{equation*}
These local ideals glue into a quasi-coherent sheaf of ideals $\mathcal{J}_i$ on $X_E$.

\end{lemma}

\begin{proof}

Equation~\eqref{eq:2.5} is the identity obtained by pulling back along $s$ the fact that
the generators of $K_W$ become zero in $B_W$. The left-hand side is tied to the residual
because, by condition~1, $u_W$ represents the evaluation, and by condition~2, $r_{W,i,a}$
represents the evaluation of the residual. The composition rule for evaluations follows
from that for ring homomorphisms. An element of the localized ideal is a finite sum of
generators written with a denominator, which gives \eqref{eq:2.6}. On each chart we take
the sheaf $\widetilde{J_i(W)}$ associated with the module $J_i(W)$ and glue along the
isomorphisms on the overlaps, which preserve the same generators. Quasi-coherence means
precisely that on each affine open set the sheaf arises in this form from a module.

\end{proof}

By the condition preserving covers and overlaps, $W\mapsto\Gamma(X_W,\mathcal{O}_{X_E})$
is a sheaf with respect to $J$. Hence $\eta$ factors through the sheafification, and we can
compare $\mathcal{I}_i^E$ on the site with $\mathcal{J}_i$ on the scheme. Mapping the local
generators of the former to the charts and taking sums with ring coefficients yields the
$J_i(W)$ of the latter. This comparison is determined by the original generators and
localization.

\begin{definition}[Lawful locus]\label{def:2.8}

Put $\mathcal{J}_{\mathrm{Ob}}=\sum_{i\in K_E^{\mathrm{req}}}\mathcal{J}_i$, and write
$X_E^{\mathrm{law}}=V(\mathcal{J}_{\mathrm{Ob}})$ for the closed subscheme that it defines.
On each chart it is $\mathrm{Spec}(B_W/\sum_iJ_i(W))$. Further,
$s^{*}_{\mathrm{ideal}}\mathcal{J}$ denotes the sheaf of ideals generated inside
$\mathcal{O}_T$ by the image of $\mathcal{J}$; it is the image of the morphism from the
module pullback $s^*\mathcal{J}$ to $\mathcal{O}_T$.

On the ring side, the two-stage conditions correspond to the two quotients
\[
D_W\longrightarrow B_W=D_W/K_W
\longrightarrow B_W/\sum_{i\in K_E^{\mathrm{req}}}J_i(W)
\]
The first quotient equates the evaluation of the equations with that of the residuals, and
the second imposes the validity of the required Laws. This order is what makes the
zero-locus condition of the next theorem test the residuals of objects.

\end{definition}

\begin{theorem}[Correspondence of equations, ideals, and zero loci]\label{thm:2.9}

Under the conditions of Construction~\ref{cons:2.6}, for every $s:T\to X_E$ the following
are equivalent.
\begin{equation*}
\begin{aligned}
&\forall W\ \forall i\in K_E^{\mathrm{req}}\ \forall a,
\quad\varepsilon_{W,i,a}(s)=0;\\
&s^{*}_{\mathrm{ideal}}\mathcal{J}_{\mathrm{Ob}}=0;\\
&s\text{ factors through }X_E^{\mathrm{law}}\text{.}
\end{aligned}
\tag{2.7}\label{eq:2.7}
\end{equation*}

\end{theorem}

\begin{proof}

That the pulled-back ideal is zero is equivalent to all of its generators becoming zero on
$T_W$. By \eqref{eq:2.5}, the vanishing of these generators coincides with the vanishing of
the residuals. Pulling back a sum of ideals and taking the generated ideal is the same as taking the sum of
the pulled-back generators, so the first two conditions are equivalent over all required
indices.

The last condition is obtained, on each affine open set, from the universal property of
homomorphisms factoring through the quotient ring. The local factorizations are unique, so
they agree on overlaps and glue. Conversely, if a factorization exists, the pullbacks of
all generators that become zero in the quotient are zero.

\end{proof}

This factorization is also the standard universal property of closed subschemes
(\cite[\href{https://stacks.math.columbia.edu/tag/01HP}{Lemma~26.4.6, Tag~01HP}]{Stacks}).
We write $\mathrm{EquationLawful}_E(s)$ for the first condition of the theorem.

\begin{example}[Synchronization of two values]\label{ex:2.10}

Let the coordinate ring be $k[x,y]$, and express the Law equating the two values by
$f=y-x$. The residual under the evaluation $x\mapsto u$, $y\mapsto v$ is $v-u$. The ring of
the lawful locus is $k[x,y]/(y-x)\cong k[x]$, and a point of it amounts to assigning one
value to both sides. In this case \eqref{eq:2.5} is verified by the substitution
$f(u,v)=v-u$.

\end{example}

\begin{corollary}[Comparison with lawfulness of objects]\label{cor:2.11}

Suppose that there is a realization $s_A:T_A\to X_E$ of an object $A$, together with, in
each context, a ring isomorphism $\Gamma((T_A)_W,\mathcal{O}_{T_A})\cong O_E(W)$ carrying
the evaluation of the residuals back to the original $\varepsilon_{W,A,i,a}$. Then
$\mathrm{Lawful}_E(A)$ of Definition~\ref{def:1.13} is equivalent to $s_A$ factoring
through $X_E^{\mathrm{law}}$.

\end{corollary}

\begin{proof}

The isomorphisms preserve and reflect zero, so the vanishing of each residual of the
object coincides with the vanishing after evaluation. Apply Theorem~\ref{thm:2.9}.

\end{proof}

\section{Coefficients representing obstructions}\label{sec:2.3}

Adding and subtracting the discrepancies of local states requires an abelian group in
which their values live. We define a way to build quotient groups from the equations and a
way to map finite failure patterns into the coefficients. The choice of coefficients
determines which differences are treated as the same.

\begin{construction}[Quotient coefficients generated from the equations]\label{cons:2.12}

In each context, put
\begin{equation*}
Q_E(W)=O_E(W)/I_{\mathrm{Ob}}^E(W),
\qquad
\overline\varepsilon_{W,A,i,a}=[\varepsilon_{W,A,i,a}]
\tag{2.8}\label{eq:2.8}
\end{equation*}
By Proposition~\ref{prop:2.4} the restrictions descend to the quotients, and $Q_E$ becomes
a presheaf with values in additive groups. The restriction of a residual class equals the
residual of the same object, Law, and Atom, restricted and then passed to the quotient.
Moreover
\[
\overline\varepsilon_{W,A,i,a}=0
\quad\Longleftrightarrow\quad
\varepsilon_{W,A,i,a}\in I_{\mathrm{Ob}}^E(W).
\]
This is the definition of the zero element of a quotient group. If the Law holds and the
residual itself is zero, then the residual class is also zero. Detection in the converse
direction requires that a failing residual not lie in the ideal.

\end{construction}

\begin{example}[Residuals distinguished in the quotient]\label{ex:2.13}

If $O_E(W)=\mathbb{Z}[x]$ and $I_{\mathrm{Ob}}^E(W)=(x)$, then $Q_E(W)\cong\mathbb{Z}$.
The class of the symbolic generator $x$ is zero, while the class of $1$, chosen as the
residual of some object, is nonzero. If the residual is $x$, it is itself nonzero but
becomes zero in the quotient. A verdict in the quotient coefficients therefore reads the
information of the specified residual modulo the ideal.

\end{example}

\begin{definition}[Local circuits and realization in coefficients]\label{def:2.14}

For an object $A$ and a context $W$, we define the Law $i$ to hold on $W$ by
\[
E_{i,W}(A)\quad\Longleftrightarrow\quad
\forall a\in\mathrm{At},\quad\varepsilon_{W,A,i,a}=0
\]
Its failure means that at least one residual is nonzero. The condition $E_i(A)$ of
Chapter~\ref{chap:1} imposes this local validity condition on all contexts.

We choose a configuration read in each context and a detector code of the form of
Definition~\ref{def:1.15}. Let $T_{i,W}$ be its finite accepted set, and write
$\mathrm{Matches}_W(Q,A)$ for the matching condition obtained by evaluating each query of
Definition~\ref{def:1.14} on this local configuration. Among the finite patterns that are
accepted and match, we select those to be used as
\[
\mathrm{Circ}^{\mathrm{loc}}_E(A,i;W)
\subseteq\{Q\in T_{i,W}\mid\mathrm{Matches}_W(Q,A)\}
\]
and call its elements local circuits. We equip this family with maps that preserve
matching and acceptance along the restrictions of the adopted contexts, and impose the
identity and composition laws. Choosing a family that remains usable after restriction in
this way is also part of the data of local detection. Local soundness, and local
completeness for the required Laws, are the conditions
\[
\begin{aligned}
c\in\mathrm{Circ}^{\mathrm{loc}}_E(A,i;W)
&\Longrightarrow\neg E_{i,W}(A),\\
i\in K_E^{\mathrm{req}}\ \land\ \neg E_{i,W}(A)
&\Longrightarrow\mathrm{Circ}^{\mathrm{loc}}_E(A,i;W)\ne\varnothing
\end{aligned}
\]
respectively. We verify these conditions for the objects, contexts, and Laws to which they
are applied.

We further choose a sheaf of abelian groups $F$. To each local circuit $c$ we associate a
witness index $a(c)\in\mathrm{At}$ and assign the generator
$\nu(c)=\nu_{W,i,a(c)}\in I_i^E(W)$. When, in addition, an additive map commuting with
restriction
\begin{equation*}
\rho_{i,W}:I_i^E(W)\longrightarrow F(W),
\qquad \kappa_{i,W}(c)=\rho_{i,W}(\nu(c))
\tag{2.9}\label{eq:2.9}
\end{equation*}
has been constructed, we call it a realization of circuits in the coefficients. The
assignment of generators is also required to commute with the restrictions of local
circuits and of the rings. For example, the sheafified quotient map
$I_i^E\to a_J(I_i^E/(I_i^E)^2)$ is a candidate for such a map. Whether its image detects
the circuits is checked on the chosen generators.

\end{definition}

\begin{proposition}[Detection of finite failures by coefficients]\label{prop:2.15}

Fix an object $A$ and a context $W$. Suppose that local soundness and local completeness
of Definition~\ref{def:2.14} hold for the required Laws, and that the image under
\eqref{eq:2.9} of each local circuit is nonzero. Then at least one required Law fails on
$W$ if and only if there exists a nonzero circuit image for some required Law.

\end{proposition}

\begin{proof}

From a failure, completeness yields a circuit, whose image is nonzero by assumption.
Conversely, if a circuit exists, soundness implies that the corresponding Law fails.

\end{proof}

When several images are aggregated into one, we also need the condition that nonzero
images do not cancel in the sum. For example, if each image is stored in a direct sum with
independent coordinates, the vanishing of the whole coincides with the vanishing of each
coordinate. The coefficients chosen at this stage are also used in the local comparisons
that follow.

\section{Comparison on overlaps and \v{C}ech obstruction classes}\label{sec:2.4}

When comparing differences of local states, we place the differences on double overlaps
and the coherence conditions among the differences on triple overlaps. The \v{C}ech
complex assembles this computation from restriction maps and addition and subtraction.

\begin{definition}[Covers and the \v{C}ech complex]\label{def:2.16}

Fix a finite cover $\mathcal{U}=(U_\alpha\to W)_{\alpha\in A}$. We take each morphism to
be injective and choose a total order on $A$. The overlaps $U_{\alpha\beta}$ and
$U_{\alpha\beta\gamma}$ are pullbacks over $W$. The coefficients $F$ form a presheaf with
values in abelian groups. On empty overlaps we set $F(\varnothing)=0$ and omit those
components from the following products.
\begin{equation*}
\begin{aligned}
C^0(\mathcal{U},F)&=\prod_\alpha F(U_\alpha),\\
C^1(\mathcal{U},F)&=\prod_{\alpha<\beta}F(U_{\alpha\beta}),\\
C^2(\mathcal{U},F)&=\prod_{\alpha<\beta<\gamma}F(U_{\alpha\beta\gamma}).
\end{aligned}
\tag{2.10}\label{eq:2.10}
\end{equation*}
We call these elements cochains. Correction amounts for local states live in degree 0,
and differences on overlaps in degree 1. We define the differentials by
\begin{equation*}
\begin{aligned}
(d^0b)_{\alpha\beta}&=b_\beta|_{U_{\alpha\beta}}-b_\alpha|_{U_{\alpha\beta}},\\
(d^1g)_{\alpha\beta\gamma}&=g_{\beta\gamma}|_{U_{\alpha\beta\gamma}}
-g_{\alpha\gamma}|_{U_{\alpha\beta\gamma}}
+g_{\alpha\beta}|_{U_{\alpha\beta\gamma}}
\end{aligned}
\tag{2.11}\label{eq:2.11}
\end{equation*}
From now on, restriction to the same overlap is abbreviated by $|$. In higher degrees,
likewise, degree $n$ consists of the product of sections on $n+1$-fold intersections,
and the differential is the alternating sum of restrictions.

\end{definition}

\begin{lemma}[Composition of the differentials]\label{lem:2.17}

We have $d^1d^0=0$.

\end{lemma}

\begin{proof}

The component restricted to a triple overlap is
$(b_\gamma-b_\beta)-(b_\gamma-b_\alpha)+(b_\beta-b_\alpha)=0$. That the two ways of
restricting each $b$ agree uses the composition rule of the presheaf.

\end{proof}

\begin{definition}[Obstruction groups on a fixed cover]\label{def:2.18}

A $g$ satisfying $d^1g=0$ is called a 1-cocycle, and an element of the form $d^0b$ a
1-coboundary. The first \v{C}ech cohomology with respect to the cover $\mathcal{U}$ is
\begin{equation*}
\check H^1(\mathcal{U},F)=\ker d^1/\mathrm{im}\,d^0
\tag{2.12}\label{eq:2.12}
\end{equation*}
Here $[g]$ denotes the equivalence class of $g$. Two cocycles have the same class when
they differ by some $d^0b$. In particular, $[g]=0$ means that there exists a local
correction $b$ with $g=d^0b$. Throughout this chapter, $\check H^1$ denotes this group
with respect to the cover and coefficients made explicit.

\end{definition}

\begin{example}[Three local references]\label{ex:2.19}

Suppose that the double overlaps of three charts are all nonempty and the triple overlap
is empty. Take the same abelian group $M$ as the coefficients on each chart and overlap,
with identity maps as restrictions. Then
\begin{equation*}
\begin{aligned}
d^0(b_0,b_1,b_2)&=(b_1-b_0,\ b_2-b_0,\ b_2-b_1),\\
\mathrm{per}(g)&=g_{01}+g_{12}-g_{02}
\end{aligned}
\tag{2.13}\label{eq:2.13}
\end{equation*}
We call the sum around the cycle, $\mathrm{per}(g)$, the period. Since there is no triple
overlap, every 1-cochain is a cocycle. The period of $d^0b$ is zero; conversely, if the
period is zero, then putting $b_0=0$, $b_1=g_{01}$, $b_2=g_{02}$ gives $d^0b=g$. Hence
$\check H^1(\mathcal{U},M)\cong M$.

If $M=\mathbb{Z}$, $g_{01}=1$, and $g_{02}=g_{12}=0$, the period is 1. No matter how
the pairwise differences of references are corrected, this difference around the cycle
remains. An actual cover of this shape is constructed from a finite topological space in
\S\ref{sec:2.8}.

\end{example}

\section{From vanishing of the obstruction class to gluing}\label{sec:2.5}

Equation~\eqref{eq:2.12} is an algebraic condition for finding correction amounts. When
the corrections act on the local states and the corrected states glue, a global state is
obtained. We describe these two conditions in terms of torsors and sheaves.

\begin{definition}[Local states and torsors]\label{def:2.20}

Take a presheaf $P$ of local states and a coefficient presheaf $F$. In this section, the
``intersections of the cover'' include each chart itself. Suppose that $F(V)$ acts on
$P(V)$ on each intersection of the cover, and that these actions commute with the
restrictions between intersections. When, on each nonempty intersection $V$ of the cover,
for any $p,q\in P(V)$ there exists a unique $m\in F(V)$ with $q=m+p$, we call the nonempty
$P(V)$ an $F(V)$-torsor.

A torsor has no designated zero state serving as a reference. Instead, the difference
$q-p=m$ of two states is uniquely determined. Every element of the coefficient group acts
as an actual change of state, so a correction amount found as a cochain can be applied to
the states. Over an empty intersection, we take $P$ to have at most one element. A family of local
states $p_\alpha\in P(U_\alpha)$ is called a local atlas.

Each $p_\alpha$ is given as an element of the state space $P(U_\alpha)$, and no
identification with the coefficient group $F(U_\alpha)$ has been chosen yet. To treat this
family as a 0-cochain of coefficients, we need a reference expressing each state in the
coefficients. Choosing a reference state locally yields such an expression, but these
references may themselves disagree on the overlaps. Whether a common reference compatible
with the restrictions can be chosen is the gluing problem studied here. In \S\ref{sec:2.8}
we make the coordinates in the local references explicit as $p$ and the shifts between
references as $\xi$, and compute the difference of states as $\xi+d^0p$.

\end{definition}

\begin{proposition}[The obstruction class defined by differences of states]\label{prop:2.21}

Defining $g_{\alpha\beta}=p_\beta| -p_\alpha|$ from a local atlas, $g$ is a cocycle.
Changing the atlas to $p'_\alpha=b_\alpha+p_\alpha$ gives
\begin{equation*}
g'=g+d^0b,
\qquad [g']=[g]
\tag{2.14}\label{eq:2.14}
\end{equation*}
We call this $[g]$ the obstruction class for gluing the local states.

\end{proposition}

\begin{proof}

On a triple overlap, the difference that yields $p_\gamma$ directly from $p_\alpha$ equals
the difference obtained via $p_\beta$. By freeness of the action,
$g_{\alpha\gamma}=g_{\beta\gamma}+g_{\alpha\beta}$, and we get $d^1g=0$. Computing the
difference after the change gives $g'_{\alpha\beta}=g_{\alpha\beta}+b_\beta|-b_\alpha|$,
so \eqref{eq:2.14} follows.

\end{proof}

\begin{theorem}[Local corrections and global states]\label{thm:2.22}

Suppose that $P$ of Definition~\ref{def:2.20} is a sheaf on the site $(\mathcal{C},J)$.
Then, for the obstruction class of a local atlas,
\begin{equation*}
P(W)\ne\varnothing
\quad\Longleftrightarrow\quad
[g]=0\text{ in }\check H^1(\mathcal{U},F)
\tag{2.15}\label{eq:2.15}
\end{equation*}
holds. If $g=d^0b$, the corrected family $(-b_\alpha)+p_\alpha$ glues uniquely into a
global state.

\end{theorem}

\begin{proof}

Suppose $g=d^0b$. By Proposition~\ref{prop:2.21} the differences after correction are
zero. On overlaps with the order reversed, the difference merely changes sign and is again
zero. On self-intersections the two projections agree by injectivity of the covering
morphisms, and on empty intersections they agree since there is at most one element. Hence
the corrected family is a matching family agreeing on all overlaps. By the sheaf condition
it glues into a unique global state.

Conversely, given $p\in P(W)$, on each chart we can take a difference with
$p_\alpha=b_\alpha+p|_{U_\alpha}$. Computing on the overlaps gives $g=d^0b$.

\end{proof}

What is unique is the gluing of the fixed corrected family. Different choices of
correction may yield different global states. For torsors under a sheaf of abelian groups,
this argument amounts to the standard criterion for trivialization
(\cite[\href{https://stacks.math.columbia.edu/tag/03AG}{\S21.4, Tag~03AG}]{Stacks}).

\subsection*{Application to local states satisfying the Laws}

\begin{corollary}[Gluing of lawful states]\label{cor:2.23}

Suppose there are a sheaf of states $\mathcal{T}$ and a natural family of required
residuals with values in the sheaf of rings $\mathcal{O}$, and let
$P(V)\subseteq\mathcal{T}(V)$ be the set of states on which all required residuals are
zero. Suppose that the action of $F$ preserves $P$ and satisfies the conditions of
Definition~\ref{def:2.20} on the chosen cover. Then the obstruction class of the local
lawful states is zero if and only if $P(W)\ne\varnothing$.

\end{corollary}

\begin{proof}

$P$ is a sheaf. Indeed, a matching family glues uniquely in $\mathcal{T}$, and each of its
residuals is zero on the cover. By separatedness of $\mathcal{O}$ it is zero globally as
well. Apply Theorem~\ref{thm:2.22}.

\end{proof}

Using the evaluation of Construction~\ref{cons:2.6}, this residual condition coincides
with factorization through the lawful locus of Theorem~\ref{thm:2.9}. To return to
lawfulness of objects, we also use the comparison of evaluations of
Corollary~\ref{cor:2.11}.

When local lawfulness is verified from finite checks, the connection between the checks
and the residuals must be specified. Concretely, we verify the following conditions.

\begin{itemize}
\item \textbf{The information can be read.} The cover is
$(E,\mathcal{R},\mathrm{Ov})$-adequate in the sense of Chapter~\ref{chap:1}, and the
necessary coordinates and interactions can be read.
\item \textbf{Failures can be detected.} Local soundness and local completeness of
Definition~\ref{def:2.14} hold, and all necessary witnesses are checked within the cover.
\item \textbf{Zero verdicts correspond.} The zero verdict on the adopted axes is
equivalent to the vanishing of the corresponding residuals.
\item \textbf{Aggregation loses no failure.} When circuit images are aggregated, each
image is nonzero and they do not cancel in the sum.
\end{itemize}

These conditions tie the local checks to membership in $P(U_\alpha)$, and the torsor
action and the sheaf condition give the global conclusion.

\section{Shape of covers, capacity of obstructions, and detection by cycles}\label{sec:2.6}

The size of the obstruction group expresses how many independent discrepancies it can
record. Whether the local states at hand carry an obstruction, on the other hand, is
determined by examining the class built from those states. In this section we study the
former by dimensions and the latter by sums along cycles.

\begin{proposition}[Dimension of the obstruction space and alternating sums]\label{prop:2.24}

Let $k$ be a field and suppose that $C^0,C^1,C^2$ are finite-dimensional $k$-vector
spaces. Then
\begin{equation*}
\dim_k\check H^1
\geq\dim_k C^1-\dim_k C^0-\dim_k C^2.
\tag{2.16}\label{eq:2.16}
\end{equation*}
Moreover, for a finite-dimensional complex of finite length,
\begin{equation*}
\chi(C^\bullet)
=\sum_n(-1)^n\dim_k C^n
=\sum_n(-1)^n\dim_k H^n(C^\bullet).
\tag{2.17}\label{eq:2.17}
\end{equation*}

\end{proposition}

\begin{proof}

Apply the rank-nullity theorem to
$\dim\check H^1=\dim\ker d^1-\dim\mathrm{im}\,d^0$. From
$\mathrm{rank}\,d^1\leq\dim C^2$ and $\mathrm{rank}\,d^0\leq\dim C^0$ we get
\eqref{eq:2.16}. Equation~\eqref{eq:2.17} follows by taking the alternating sum of
$\dim C^n=\dim\mathrm{im}\,d^{n-1}+\dim H^n+\dim\mathrm{im}\,d^n$, in which the dimensions
of the images cancel between adjacent degrees.

\end{proof}

We call \eqref{eq:2.16} the lower bound on obstruction capacity. A change preserving the
cochain dimensions in each degree preserves $\chi$, but the dimensions of the individual
cohomology groups and a specified class $[g]$ also involve the differentials and $g$
itself.

\begin{proposition}[Constant coefficients and the nerve of a cover]\label{prop:2.25}

Take a finite open cover of a topological space and suppose that all of its nonempty
finite intersections are connected. Take as coefficients the sheaf of locally constant
$k$-valued functions. Let $N(\mathcal{U})$ be the complex with a simplex for each
nonempty intersection $U_{\alpha_0\cdots\alpha_n}$. Then the \v{C}ech complex is
isomorphic to the usual cochain complex of $N(\mathcal{U})$ with coefficients in $k$, and
$\dim_k\check H^1(\mathcal{U},k)=\beta_1(N(\mathcal{U});k)$.

\end{proposition}

\begin{proof}

A locally constant function on a connected intersection is determined by a single value.
Under this identification the restrictions become identity maps, and \eqref{eq:2.11}
coincides with the alternating sum along the faces of a simplex. Taking this
identification in each degree gives the isomorphism of complexes. Moreover,
$\beta_1(N(\mathcal{U});k)$ is the dimension of the first cohomology of this complex.

\end{proof}

When an intersection has several connected components, we place one coordinate per
component. When the components split into finitely many open connected parts, the same
construction can be carried out on the complex whose vertices, edges, and faces are the
components.

\begin{proposition}[Vanishing on forests]\label{prop:2.26}

Suppose that there are no triple overlaps and that the finite graph whose vertices are the
charts and whose edges are the components of the nonempty overlaps is a forest. If the
coefficients decompose along the components and the restriction
$F(U_\alpha)\to F(U_e)$ from each endpoint to each edge component is surjective, then
$\check H^1(\mathcal{U},F)=0$.

\end{proposition}

\begin{proof}

Choose a root in each tree and set the correction at the root to zero. Once the correction
of a parent is determined, choose the correction of each child so that the difference
along the parent-child edge equals the specified $g_e$. Surjectivity of the restriction to
the edges makes this choice possible. Since each child has only one parent, we can proceed
without altering the conditions on edges already satisfied. After processing all vertices
we have $d^0b=g$.

\end{proof}

\begin{example}[Paths and cycles]\label{ex:2.27}

On a three-vertex path with identity restrictions, the two differences can be absorbed in
order from the root. Adding an edge joining the two ends leaves the period of
Example~\ref{ex:2.19}. If moreover there is a triple overlap with the same coefficients
and restrictions, the cocycle condition of \eqref{eq:2.11} makes this period zero. Thus
not only the edges of the overlaps but also the faces filling the space between them enter
the computation of the obstruction group.

\end{example}

\subsection*{Sums along cycles and differences on the boundary}

\begin{proposition}[Periods and the Stokes formula]\label{prop:2.28}

Pair the chains and cochains of a finite complex $N$ with constant coefficients in $k$ by
multiplying the coefficients of simplices with matching orientation and summing. For the
boundary map $\partial$ of chains and the differential $d$ of cochains,
\begin{equation*}
\langle d\omega,z\rangle=\langle\omega,\partial z\rangle
\tag{2.18}\label{eq:2.18}
\end{equation*}
holds. Hence $\langle g,z\rangle$ over a 1-cycle $z$ is unchanged when a coboundary is
added to $g$. If this value is nonzero, then $[g]\ne0$.

\end{proposition}

\begin{proof}

For a single simplex, both sides are the sum of the values on its faces with the same
alternating signs. Linearity extends this to general chains. If $\partial z=0$, then
$\langle d^0b,z\rangle=0$, and the remaining claims follow.

\end{proof}

\begin{proposition}[Obstructions arising from differences on the boundary]\label{prop:2.29}

Split a finite complex into two subcomplexes, $N=N_A\cup N_B$, and put $L=N_A\cap N_B$.
For 0-cochains with constant coefficients, the difference of restrictions
$(u,v)\mapsto u|_L-v|_L$ is surjective. Take $b\in H^0(L,k)$, that is, a 0-cochain with
$db=0$, and choose 0-cochains $u,v$ with $u|_L-v|_L=b$. Then $du,dv$ glue into a
single 1-cocycle $h$, and $\delta(b)=[h]\in H^1(N,k)$ does not depend on the choices.
Moreover,
\begin{equation*}
\begin{aligned}
\delta(b)=0\quad\Longleftrightarrow\quad
&\exists a\in H^0(N_A,k),\ \exists c\in H^0(N_B,k),\\
&b=a|_L-c|_L.
\end{aligned}
\tag{2.19}\label{eq:2.19}
\end{equation*}

\end{proposition}

\begin{proof}

Surjectivity is obtained by extending the value on each simplex of $L$ to $N_A$. Since
$d(u|_L-v|_L)=db=0$, the cochains $du,dv$ agree and glue into $h$. That $dh=0$
follows from $d^2=0$ on both parts. The difference from another lift is the restriction of
a 0-cochain $w$ on $N$, and $h$ changes by $dw$. Moreover, if $h=dw$, then
$a=u-w|_{N_A}$ and $c=v-w|_{N_B}$ are 0-cocycles whose difference is $b$. For the
converse, it suffices to choose $a,c$ as the lifts.

\end{proof}

We can read this $\delta(b)$ as the obstruction created by the difference $b$ on the
boundary. Splitting a cycle $z=z_A+z_B$ into chains on the two parts,
$\partial z_A=-\partial z_B$ is supported on $L$. By Proposition~\ref{prop:2.28},
$\langle h,z\rangle=\langle b,\partial z_A\rangle$. The difference around the whole cycle
coincides with the signed sum of the differences compared on the boundary of the two
parts.

\section{Comparison of the structural and semantic phases}\label{sec:2.7}

Whether a fact belongs to structure or to semantics is determined by what changes under
the permitted changes of semantics. From this classification we split the coefficients and
study conditions under which no obstruction is lost by looking only at the semantic side.

\begin{definition}[Two phases by invariance of extraction]\label{def:2.30}

For the extraction doctrine $D$ of Chapter~\ref{chap:1}, choose a family
$(D_\lambda)_{\lambda\in\Lambda}$ that preserves the vocabulary, the normalization, the
resolution, and the structure of sources, and changes only the semantic reading and its
admissible predicates. We require this family to contain $D$ itself. A pair $(s,a)$ of a
source and an Atom is structural if
\begin{equation*}
\forall\lambda\in\Lambda,\quad
\bigl(a\in\mathrm{Atomize}_{D_\lambda}(s)
\Longleftrightarrow a\in\mathrm{Atomize}_D(s)\bigr)
\tag{2.20}\label{eq:2.20}
\end{equation*}
holds. A pair that is not structural is called semantic. For example, if the existence of
a component is extracted under every reading while only a fact about the interpretation
of values ceases to be extracted under some reading, then the former is structural and
the latter semantic.

\end{definition}

\begin{definition}[The two-phase coefficient complex]\label{def:2.31}

Assign to each coordinate of a finite-dimensional complex $C^\bullet_{\mathrm{all}}$ a
pair of a source and an Atom. Let $C^n_{\mathrm{str}}$ be the subspace spanned by the
coordinates belonging to structural pairs. We write Condition~S for the requirement that the
restrictions and the differentials preserve this part. Under this condition the
differential descends to $C^n_{\mathrm{sem}}=C^n_{\mathrm{all}}/C^n_{\mathrm{str}}$, and
\begin{equation*}
0\longrightarrow C^\bullet_{\mathrm{str}}
\longrightarrow C^\bullet_{\mathrm{all}}
\longrightarrow C^\bullet_{\mathrm{sem}}\longrightarrow0
\tag{2.21}\label{eq:2.21}
\end{equation*}
is an exact sequence in each degree. This $C^\bullet_{\mathrm{sem}}$ is the quotient
complex built from the two phases of extraction. In \S\ref{sec:2.9} we will separately
build coefficients $M_{\mathrm{sem}}$ from the relations of repair operations.

\end{definition}

\begin{proposition}[Retention of obstructions on the semantic side]\label{prop:2.32}

If Condition~S holds and $H^1(C^\bullet_{\mathrm{str}})=0$, then the induced map
\begin{equation*}
H^1(C^\bullet_{\mathrm{all}})
\longrightarrow H^1(C^\bullet_{\mathrm{sem}})
\tag{2.22}\label{eq:2.22}
\end{equation*}
is injective.

\end{proposition}

\begin{proof}

Suppose that a cocycle $g$ becomes a coboundary on the semantic side. Lift the
corresponding semantic 0-cochain to some $b\in C^0_{\mathrm{all}}$ by surjectivity in
degree 0. Then $g-d^0b$ is a cocycle belonging to the structural part. Since $H^1=0$ on
the structural side, there is a structural $c$ with $g-d^0b=d^0c$. Hence $g=d^0(b+c)$.

\end{proof}

If the structural complex satisfies the forest condition of Proposition~\ref{prop:2.26},
the required $H^1=0$ is obtained. The roles of the two conditions are visible in the
following small examples.

\begin{example}[A case where Condition S fails]\label{ex:2.33}

Let $s$ be a pair that is always extracted before and after the semantic changes, and $t$
a pair whose extraction truth value changes. On the complex with two vertices and one
edge, assign $s$ to the left vertex and $t$ to the right vertex and the edge. Then
$C^0_{\mathrm{all}}=k^2$, $C^1_{\mathrm{all}}=k$, and $d^0(a,b)=b-a$; the structural part
is $k(1,0)$ in degree 0 and zero in degree 1. Since $d^0(1,0)=-1$, the structural
part is not preserved, and this differential does not descend to the semantic quotient.

\end{example}

\begin{example}[A case where Condition S alone does not give injectivity]\label{ex:2.34}

Take two copies of the complex joining two vertices by two parallel edges, and assign $s$
to all coordinates of one copy and $t$ to all coordinates of the other. The differential
of each copy is $(a,b)\mapsto(b-a,b-a)$, and degree 2 is zero. The structural part is
exactly the first copy, so Condition~S holds. However, the 1-cochain $(1,0)$ of the
first copy is not a coboundary, because the values on the two edges differ. This nonzero
class maps to zero on the semantic side. It is thus the condition that obstructions
vanish on the structural side that supports the injectivity of
Proposition~\ref{prop:2.32}.

\end{example}

\section{Concrete integer-coefficient obstructions built from Atoms}\label{sec:2.8}

To compare with the diagnosis in the next chapter, we construct coefficients and covers
concretely. The generators of the coefficients are pairs of a Law and a source, and the
points of the geometry represent the ranges read locally. We keep both in a single Atom
family and compute obstruction classes from local states and transitions.

\subsection*{Building coefficients from generator relations}

\begin{definition}[Generators and primitive relations]\label{def:2.35}

Take a finite set of sources $S$ and a finite family of Laws $(\ell)_{\ell\in L}$. Each
Law has a set of values $\mathrm{Val}_\ell$ and an evaluation
$v_\ell:S\to\mathrm{Val}_\ell$. Let the set of generators be $G=L\times S$, and let the
label of $g=(\ell,s)$ be $(\ell,v_\ell(s))$. A primitive relation $R\subseteq G\times G$
is declared as a relation preserving these labels. Let $\mathfrak{B}$ be the quotient set
by the reflexive, symmetric, and transitive closure of $R$, and write
$[g]_R\in\mathfrak{B}$ for the component of a generator $g$. We define the
integer-coefficient group by
\begin{equation*}
M_R=\mathbb{Z}^{(G)}/\langle e_g-e_h\mid gRh\rangle
\tag{2.23}\label{eq:2.23}
\end{equation*}
Here $\mathbb{Z}^{(G)}$ is the free abelian group with basis $e_g$.

\end{definition}

\begin{lemma}[Standard form of the coefficients]\label{lem:2.36}

The map $e_g\mapsto e_{[g]_R}$ induces an isomorphism
$M_R\cong\mathbb{Z}^{(\mathfrak{B})}$.

\end{lemma}

\begin{proof}

Basis elements joined by a primitive relation map to the same component, so the map
descends to the quotient. In the other direction, choose a representative generator of
each component and assign its class in $M_R$. Two generators in the same component are
joined by a finite sequence of primitive relations, and their difference is the sum of
the differences of these relations, so the assignment does not depend on the choice. The
composites of the two maps are the identity on the bases, so they are mutually inverse.

\end{proof}

As a concrete example, let $S=\{0,1\}^2$, $L=\{\ell\}$, and $v_\ell(v,h)=v$. The
primitive relation consists of the two pairs joining $(v,0)$ and $(v,1)$ with the same
$v$. Writing the generators as $g_{00},g_{01},g_{10},g_{11}$, there are two components
$B_0=\{g_{00},g_{01}\}$ and $B_1=\{g_{10},g_{11}\}$, and $M_R\cong\mathbb{Z}^2$. Even
though the second component of the original sources is identified, the two components
with different Law values remain distinguished.

\subsection*{The eight-point space and two covers}

\begin{construction}[Atom input carrying points and generators]\label{cons:2.37}

Take a set $X$ with four vertices $a_0,a_1,b,c$ and four edge points $e_{01},ab,bc,ac$.
The endpoints of the edge points are $e_{01}:(a_0,a_1)$, $ab:(a_1,b)$, $bc:(b,c)$, and
$ac:(a_0,c)$. The order is defined only by the relations placing each vertex below the
edge points incident to it, together with reflexivity. The open sets are the upward
closed subsets.

We attach tags distinguishing point Atoms from generator Atoms, and let the Atom family
be $X\sqcup G$. A point Atom has the point itself as its subject, with fixed values for
the predicate, the payload, and the axis. A generator Atom $(\ell,s)$ has the source $s$,
the Law index $\ell$, and the Law value $v_\ell(s)$ as its subject, predicate, and
payload, respectively, with a single fixed axis. Each Atom is uniquely determined by its
kind and these coordinates. The relation of the configuration is the $R$ declared between
the generator Atoms.

The context $c_V$ corresponding to an open set $V\subseteq X$ reads the point Atoms of
$V$ and all generator Atoms. Inclusions $V'\subseteq V$ are the restriction morphisms,
and the overlaps are given by $c_{V'\cap V}$. For covers we adopt families that cover the
open set and also cover the generators that are read. We use the Grothendieck topology on
the context site generated by these families. In a nonempty covering family each context
reads all generators, so the covering condition on generators also holds. The two covers
below are covers of this context site.

\end{construction}

\begin{lemma}[The three-chart and four-chart covers]\label{lem:2.38}

The minimal open neighborhoods of the vertices are
\begin{equation*}
\begin{aligned}
U_{a_0}&=\{a_0,e_{01},ac\},& U_{a_1}&=\{a_1,e_{01},ab\},\\
U_b&=\{b,ab,bc\},& U_c&=\{c,bc,ac\}.
\end{aligned}
\tag{2.24}\label{eq:2.24}
\end{equation*}
These cover $X$. Setting $U_a=U_{a_0}\cup U_{a_1}$, the family $(U_a,U_b,U_c)$ is also a
cover. In both covers, the charts and the nonempty double overlaps are connected, and the
intersection of three distinct charts is empty.

\end{lemma}

\begin{proof}

Each vertex lies in its own neighborhood and each edge point lies in the neighborhoods of
its two endpoints, so these families are covers. Each minimal neighborhood has a least
element. Any attempt to separate it into two relatively open sets fails, because the side
containing the least element contains all points above it. Since $U_{a_0}$ and $U_{a_1}$
meet in $e_{01}$, their union is also connected. The nonempty double intersections of the
four-chart cover are $\{e_{01}\},\{ab\},\{bc\},\{ac\}$, and those of the three-chart
cover are $\{ab\},\{bc\},\{ac\}$. Each edge point has only two endpoints, so even after
merging the charts as above, no point is shared by three of them.

\end{proof}

\begin{figure}[htbp]
\centering
% Figure 2.1 -- the eight-point space and its two covers.
% Same content as the Japanese manuscript's figure ../../figures/ch02-finite-covers.svg, redrawn in TikZ.
\begin{tikzpicture}[
    x=1mm, y=1mm,
    vertex/.style={circle, fill, inner sep=0pt, minimum size=4.6pt},
    edgepoint/.style={rectangle, draw, semithick, fill=white, inner sep=0pt, minimum size=5pt},
    order/.style={-{Stealth[length=4.5pt]}, semithick},
    nerve/.style={semithick},
    every label/.style={font=\normalsize},
    panel/.style={font=\small},
    note/.style={font=\small}
  ]
  % (a) The order on the eight-point space: circles are vertices, squares are
  % edge points, and the arrows point from a vertex to a larger edge point.
  \begin{scope}
    \node[panel, align=center] at (16,44) {(a) Order on the\\ eight-point space};
    \node[vertex, label=above:{$a_0$}] (a0) at (0,32) {};
    \node[vertex, label=above:{$a_1$}] (a1) at (32,32) {};
    \node[vertex, label=below:{$b$}]   (b)  at (32,0)  {};
    \node[vertex, label=below:{$c$}]   (c)  at (0,0)   {};
    \node[edgepoint, label=above:{$e_{01}$}] (e01) at (16,32) {};
    \node[edgepoint, label=right:{$ab$}]     (ab)  at (32,16) {};
    \node[edgepoint, label=below:{$bc$}]     (bc)  at (16,0)  {};
    \node[edgepoint, label=left:{$ac$}]      (ac)  at (0,16)  {};
    \draw[order] (a0) -- (e01);
    \draw[order] (a1) -- (e01);
    \draw[order] (a1) -- (ab);
    \draw[order] (b)  -- (ab);
    \draw[order] (b)  -- (bc);
    \draw[order] (c)  -- (bc);
    \draw[order] (c)  -- (ac);
    \draw[order] (a0) -- (ac);
    \node[note] at (16,-11) {arrows: vertex $<$ edge point};
  \end{scope}
  % (b) The nerve of the four-chart cover: four edges and no face.
  \begin{scope}[xshift=50mm]
    \node[panel, align=center] at (16,44) {(b) Nerve of the\\ four-chart cover};
    \node[vertex, label=above:{$U_{a_0}$}] (Ua0) at (0,32) {};
    \node[vertex, label=above:{$U_{a_1}$}] (Ua1) at (32,32) {};
    \node[vertex, label=below:{$U_b$}]     (Ub)  at (32,0)  {};
    \node[vertex, label=below:{$U_c$}]     (Uc)  at (0,0)   {};
    \draw[nerve] (Ua0) -- (Ua1) node[midway, above, note] {$\{e_{01}\}$};
    \draw[nerve] (Ua1) -- (Ub)  node[midway, right, note] {$\{ab\}$};
    \draw[nerve] (Ub)  -- (Uc)  node[midway, below, note] {$\{bc\}$};
    \draw[nerve] (Uc)  -- (Ua0) node[midway, left, note]  {$\{ac\}$};
    \node[note] at (16,-11) {edge labels: double intersections};
  \end{scope}
  % (c) The nerve of the three-chart cover: a triangle of edges with no face.
  \begin{scope}[xshift=100mm]
    \node[panel, align=center] at (16,44) {(c) Nerve of the\\ three-chart cover};
    \node[vertex, label=above:{$U_a$}] (Ua) at (16,32) {};
    \node[vertex, label=below:{$U_b$}] (Vb) at (32,0)  {};
    \node[vertex, label=below:{$U_c$}] (Vc) at (0,0)   {};
    \draw[nerve] (Ua) -- (Vb) node[midway, right=2pt, note] {$\{ab\}$};
    \draw[nerve] (Vb) -- (Vc) node[midway, below, note]     {$\{bc\}$};
    \draw[nerve] (Vc) -- (Ua) node[midway, left=2pt, note]  {$\{ac\}$};
    \node[note] at (16,-11) {$U_a = U_{a_0} \cup U_{a_1}$};
  \end{scope}
\end{tikzpicture}

\caption{\textbf{The eight-point space and the two covers.} The left panel shows the
eight points of the space $X$: circles are vertices, squares are edge points, and arrows
indicate the order from vertices to edge points. The middle and right panels show the
nerves of the covers, with nodes for charts and edges for nonempty double intersections.
When $U_{a_0}$ and $U_{a_1}$ in the middle are merged into $U_a$, their intersection
point $e_{01}$ moves into the interior of $U_a$. The triangle on the right has no face,
and the triple intersection is empty.}\label{fig:2.1}
\end{figure}

In this computation, the vanishing of degree 2 comes from the actual intersections
being empty.

\begin{construction}[The sheaf of locally constant obstruction coefficients]\label{cons:2.39}

Let $F_R(c_V)$ be the group of locally constant $M_R$-valued functions on $V$, with
restriction of functions as restriction. Functions agreeing on overlaps glue uniquely,
and local constancy is checked locally, so $F_R$ is a sheaf. On the empty open set it is
the zero group. The map $c_V\mapsto V$ from contexts to open sets preserves intersections
and covers, so these coefficients form an actual sheaf on the chosen context site.

By the connectedness of Lemma~\ref{lem:2.38}, a section over a chart or a nonempty
intersection is determined by a single value in $M_R$. In this presentation the
restrictions are identity homomorphisms. Hence the \v{C}ech complexes of both covers are
complexes with $M_R$-values on vertices and edges and zero in degree 2.

\end{construction}

\subsection*{Computing obstructions from local states and affine transitions}

\begin{proposition}[Lawful chart states]\label{prop:2.40}

Through the isomorphism of Lemma~\ref{lem:2.36}, read each chart state $p_\alpha\in M_R$
as componentwise integers $(p_{\alpha,B})_{B\in\mathfrak{B}}$. Define the polynomial ring
and the relation ideal by
\begin{equation*}
\begin{aligned}
P_R&=\mathbb{Z}[X_g\mid g\in G],\\
J_R&=(X_g-X_h\mid gRh),\\
e_\alpha:P_R&\longrightarrow\mathbb{Z},
\qquad X_g\longmapsto p_{\alpha,[g]_R}
\end{aligned}
\tag{2.25}\label{eq:2.25}
\end{equation*}
This evaluation sends $J_R$ to zero and gives an integer point of
$\mathrm{Spec}(P_R/J_R)$. The states on the overlaps obtained by actual restriction are
also lawful. Moreover, the states after translation by an arbitrary $\xi_e\in M_R$ are
lawful.

\end{proposition}

\begin{proof}

If $gRh$, then $[g]_R=[h]_R$, so the evaluation of each generator $X_g-X_h$ is zero.
Restriction to an overlap preserves the integer value of each component. After a
translation, the same $\xi_{e,B}$ is added to the same component, so the identities are
preserved.

\end{proof}

The polynomial equations here require the coordinates to agree along the declared
primitive relations. The Law values of Definition~\ref{def:2.35} retain the origin of
each generator, and in the next chapter we use them to map these coefficients into the
diagnosis.

\begin{proposition}[Obstruction of specified affine local data]\label{prop:2.41}

Write $\mathcal{U}_q$ for either of the covers and choose local data $x=(\xi,p)$. Here
$p\in C^0(\mathcal{U}_q,F_R)$ consists of the chart states, and
$\xi\in C^1(\mathcal{U}_q,F_R)$ is the translation carrying the state on the right into
the coordinates on the left across each overlap. The orientation of the edges is chosen
by the order of the charts. When an orientation is reversed, the signs of the transition
and of the comparison value are reversed. The comparison value on an edge
$e:\alpha\to\beta$ is
\begin{equation*}
\begin{aligned}
o_q(x)_e&=\xi_e+\mathrm{res}_{\beta,e}(p_\beta)
-\mathrm{res}_{\alpha,e}(p_\alpha),\\
o_q(x)&=\xi+d^0p.
\end{aligned}
\tag{2.26}\label{eq:2.26}
\end{equation*}
This is the difference obtained by actually restricting the lawful states on the two
sides, carrying one of them by the transition, and then comparing. Since degree 2 is
zero, it is a cocycle and defines $[o_q(x)]\in\check H^1(\mathcal{U}_q,F_R)$. Changing
the local states to $p+h$ gives
\begin{equation*}
o_q(\xi,p+h)=o_q(\xi,p)+d^0h,
\qquad [o_q(\xi,p+h)]=[o_q(\xi,p)]
\tag{2.27}\label{eq:2.27}
\end{equation*}
\end{proposition}

\begin{proof}

Equation~\eqref{eq:2.26} computes the difference between the transported right state and
the left state, and by Proposition~\ref{prop:2.40} both states being compared are lawful.
Additivity of the differential gives \eqref{eq:2.27}.

This class is also obtained as the gluing obstruction of Proposition~\ref{prop:2.21}.
Place a locally constant $M_R$-valued state on each chart and compare coordinates from
right to left on each overlap by $m\mapsto\xi_e+m$. The reverse direction uses $-\xi_e$,
and self-comparison is zero. Since the intersection of three distinct charts is empty, no
additional condition arises on triple overlaps. Defining the states over an open set as
the families of local states satisfying these comparisons yields a sheaf. On each chart
and each nonempty intersection, the action of the coefficients $F_R$ is free and
transitive, and the difference of the $p$ chosen in chart coordinates equals
\eqref{eq:2.26}. Hence the same cochain yields the same obstruction class.

\end{proof}

\begin{example}[Changing local states and changing transitions]\label{ex:2.42}

Order the three-chart cover by $a<b<c$ and let $u=e_{B_0}\in M_R\cong\mathbb{Z}^2$.
Choose $p=0$, $\xi_{ab}=u$, and $\xi_{ac}=\xi_{bc}=0$. The period of
Example~\ref{ex:2.19} is $u$, so $[o_q(x)]\ne0$. The chart states each satisfy the Laws of
Proposition~\ref{prop:2.40}, but under the fixed transitions they do not glue into a
global state.

Changing only $p$ does not change the class, by \eqref{eq:2.27}. On the other hand,
changing $\xi_{ab}$ itself to zero leaves $p=0$ compatible as it is, and the obstruction
becomes zero. This computation distinguishes the differences removable by changing local
coordinates from those removable only by changing the transitions.

The data passed to Chapter~\ref{chap:3} are the coefficients $F_R$, the cover
$\mathcal{U}_q$, the local data $x$, and the concrete class $[o_q(x)]$ generated from
them.

\end{example}

\section{Semantic local repair and obstruction classes of equations}\label{sec:2.9}

On the repair side, we build coefficients from the available changes and the relations
identifying changes with the same effect. On the equation side, we use the quotient
coefficients of \S\ref{sec:2.3}. After building an obstruction class from the local
states on each side, we compare the two in the next section.

Throughout, we fix a finite cover $\mathcal{U}$ by injective morphisms as in
Definition~\ref{def:2.16}. We call the diagram consisting of the charts, the nonempty
double and triple intersections, and their restriction morphisms the intersection
diagram of the cover. The algebraic conditions needed for the comparison are imposed on
this diagram.

\begin{definition}[Presentation of semantic repair]\label{def:2.43}

To each context $V$ we assign a set $\Lambda(V)$ of Atoms distinguished on the semantic
side, a map $\pi_V:\Lambda(V)\to\mathrm{At}$ giving the underlying Atom, and a subset
$S(V)\subseteq\Lambda(V)$ usable for repair. These carry functorial restrictions; the
restrictions preserve $S$ and commute with $\pi$. Here the underlying Atoms are read as
the same under restriction.

A repair word is a finite integer combination of elements of $S(V)$. We declare, as a subgroup
$R_{\mathrm{rep}}(V)$, the differences of words regarded as locally having the same
effect, and require this subgroup to be preserved by restriction as well.
We define the semantic coefficients by
\begin{equation*}
\begin{aligned}
F_{\mathrm{rep}}(V)&=\mathbb{Z}^{(S(V))},\\
M_{\mathrm{sem}}(V)&=F_{\mathrm{rep}}(V)/R_{\mathrm{rep}}(V)
\end{aligned}
\tag{2.28}\label{eq:2.28}
\end{equation*}
The restrictions extend linearly to the free abelian groups and preserve the relation
subgroups, so they descend to the quotients. The identity and composition rules hold on
generators, so $M_{\mathrm{sem}}$ is a presheaf with values in abelian groups.

\end{definition}

\begin{proposition}[Torsors obtained from the action of repair words]\label{prop:2.44}

Suppose that $F_{\mathrm{rep}}$ acts on a presheaf $P_{\mathrm{sem}}$ of semantic local
repair states and that the action commutes with restriction. On the intersection diagram
we impose the following three conditions.

\begin{itemize}
\item If $w\in R_{\mathrm{rep}}(V)$, then $w+p=p$ for every state.
\item If $w+p=p$, then $w\in R_{\mathrm{rep}}(V)$.
\item For any two states $p,q$, some repair word $w$ satisfies $q=w+p$.
\end{itemize}

Then each nonempty $P_{\mathrm{sem}}(V)$ is an $M_{\mathrm{sem}}(V)$-torsor.

\end{proposition}

\begin{proof}

By the first condition, the action descends to the quotient by the relations. By the
second condition, a repair that is nonzero in the quotient fixes no state, so the action
is free. The third condition gives transitivity.

\end{proof}

The first two conditions state that the repair words declared to have zero effect
coincide with the words that actually leave every state unchanged. The third condition
states that the states being compared can reach one another by the chosen repairs.

\begin{definition}[Two local states and their obstruction classes]\label{def:2.45}

On the semantic side, choose a local atlas $p_\alpha\in P_{\mathrm{sem}}(U_\alpha)$. On
the equation side, choose $Q_E$ of Construction~\ref{cons:2.12} and a presheaf $P_E$ of
local states on which it acts. We require $P_E(V)$ on each nonempty intersection to be a
$Q_E(V)$-torsor, with restrictions commuting with the action. Let its local atlas be
$e_\alpha\in P_E(U_\alpha)$. An element of $P_E$ represents a state that concretely
realizes the chosen local equation data.

With the two coefficients we form the \v{C}ech complexes
$C^\bullet_{\mathrm{rep}}=C^\bullet(\mathcal{U},M_{\mathrm{sem}})$ and
$C^\bullet_E=C^\bullet(\mathcal{U},Q_E)$, and define the residuals by
\begin{equation*}
\begin{aligned}
r_{\mathrm{sem},\alpha\beta}&=p_\beta|-p_\alpha|,\\
r_{E,\alpha\beta}&=e_\beta|-e_\alpha|
\end{aligned}
\tag{2.29}\label{eq:2.29}
\end{equation*}
By Proposition~\ref{prop:2.21}, both are cocycles and their classes do not depend on the
choice of atlas. We abbreviate the semantic group
$\check H^1(\mathcal{U},M_{\mathrm{sem}})$ to $H^1_{\mathrm{sem}}(\mathcal{U})$.

On empty intersections, we take the two coefficients to be the zero groups and the
two state sets to have at most one element. Since the local atlases restrict, the state set on each
nonempty intersection is nonempty.

\end{definition}

\begin{construction}[Correspondence between Atoms and equations]\label{cons:2.46}

To each $\lambda\in S(V)$ we assign a required equation index $i_\lambda$ and a reading
$A_\lambda$ of an object, preserved under restriction. We define the map sending a
repair generator to the residual class of the corresponding equation by
\begin{equation*}
\chi_V(\lambda)
=[\varepsilon_{V,A_\lambda,i_\lambda,\pi_V(\lambda)}]
\in Q_E(V)
\tag{2.30}\label{eq:2.30}
\end{equation*}
By the restriction rule for residuals, $\chi$ commutes with restriction. By the
universal property of free abelian groups, it extends uniquely to a homomorphism
$\widetilde\chi_V:F_{\mathrm{rep}}(V)\to Q_E(V)$. The map $\chi$ specifies which repair
generator represents the residual of which object, equation, and Atom.

\end{construction}

\section{The SAGA comparison and global repair}\label{sec:2.10}

We now connect the two constructions above by a correspondence of local states. To
obtain an isomorphism of coefficients, we need the relations among repair words to hold
on the equation side as well, every repair word lost on the equation side to be
contained in the relations, and the coefficients on the equation side to be expressible
by the repair generators. The first property can be derived from the correspondence of
states.

\begin{definition}[Correspondence of local states and completeness]\label{def:2.47}

On the intersection diagram, we construct maps
$\beta_V:P_{\mathrm{sem}}(V)\to P_E(V)$ commuting with restriction, and require, for
each generator,
\begin{equation*}
\beta_V(\lambda+p)=\chi_V(\lambda)+\beta_V(p)
\tag{2.31}\label{eq:2.31}
\end{equation*}
The $\lambda$ on the left-hand side denotes the action of a basis element of the free
abelian group. Furthermore, on each nonempty intersection $V$ we require
\begin{equation*}
\ker\widetilde\chi_V\subseteq R_{\mathrm{rep}}(V),
\qquad
\mathrm{im}\,\widetilde\chi_V=Q_E(V)
\tag{2.32}\label{eq:2.32}
\end{equation*}
We call the first condition completeness of relations and the second completeness of
generators. The former says that the repair relations account for all words whose effect
becomes zero on the equation side, and the latter that each coefficient on the equation
side can be expressed by a repair word.

\end{definition}

\begin{lemma}[The correspondence of states gives soundness of the relations]\label{lem:2.48}

Under the conditions of Proposition~\ref{prop:2.44} and
Definitions~\ref{def:2.45} and~\ref{def:2.47},
\begin{equation*}
R_{\mathrm{rep}}(V)\subseteq\ker\widetilde\chi_V
\tag{2.33}\label{eq:2.33}
\end{equation*}
holds.

\end{lemma}

\begin{proof}

Extending \eqref{eq:2.31} to finite sums and additive inverses, we get
$\beta_V(w+p)=\widetilde\chi_V(w)+\beta_V(p)$ for every repair word $w$. Let
$w\in R_{\mathrm{rep}}(V)$ and take a state $p$ obtained from the local atlas; then
\[
\beta_V(p)=\beta_V(w+p)
=\widetilde\chi_V(w)+\beta_V(p).
\]
Since the action on the equation side is free, $\widetilde\chi_V(w)=0$.

\end{proof}

\begin{theorem}[The SAGA comparison of coefficients and obstruction classes]\label{thm:2.49}

Under the data and conditions from Definition~\ref{def:2.43} through
Definition~\ref{def:2.47}, we obtain a natural isomorphism on the intersection diagram
\begin{equation*}
\Phi:M_{\mathrm{sem}}\xrightarrow{\sim}Q_E,
\qquad\Phi_V([w])=\widetilde\chi_V(w)
\tag{2.34}\label{eq:2.34}
\end{equation*}
It induces isomorphisms of the complexes in degrees 0 through 2 and of the first
cohomology, and matches the obstruction classes built from the two local atlases.
\begin{equation*}
\begin{aligned}
\kappa^n &: C^n_{\mathrm{rep}}\xrightarrow{\sim}C^n_E,
\qquad (\kappa^nc)_V=\Phi_V(c_V)\quad(n=0,1,2),\\
\kappa^{n+1}d^n_{\mathrm{rep}}&=d^n_E\kappa^n
\quad(n=0,1),
\end{aligned}
\tag{2.35}\label{eq:2.35}
\end{equation*}
\begin{equation*}
\begin{aligned}
\kappa_* &: H^1_{\mathrm{sem}}(\mathcal{U})\xrightarrow{\sim}\check H^1(\mathcal{U},Q_E),\\
\kappa_*[r_{\mathrm{sem}}]&=[r_E].
\end{aligned}
\tag{2.36}\label{eq:2.36}
\end{equation*}

\end{theorem}

\begin{proof}

We first compare the coefficients. By Lemma~\ref{lem:2.48}, $\widetilde\chi_V$ descends
to the quotient and gives $\Phi_V$. If $\Phi_V([w])=0$, then completeness of relations
gives $w\in R_{\mathrm{rep}}(V)$, so $[w]=0$. By completeness of generators, $\Phi_V$ is
also surjective. Commutation with restriction holds for $\chi$ and extends from the
generators to the whole. This gives \eqref{eq:2.34}.

Next we apply $\Phi_V$ to the components on each intersection. By additivity and
commutation with restriction, in degree 0 for example
\[
\begin{aligned}
(\kappa^1d^0_{\mathrm{rep}}b)_{\alpha\beta}
&=\Phi_{U_{\alpha\beta}}(b_\beta|-b_\alpha|)\\
&=\Phi_{U_\beta}(b_\beta)|-\Phi_{U_\alpha}(b_\alpha)|\\
&=(d^0_E\kappa^0b)_{\alpha\beta}.
\end{aligned}
\]
The same computation applies to the alternating sum of three terms in degree 1. Using
the inverse of each $\Phi_V$ componentwise gives the inverse map of complexes, so
cocycles and coboundaries are preserved in both directions. Hence
$[c]\mapsto[\kappa^1c]$ gives the isomorphism of \eqref{eq:2.36}.

Finally we compare the specified residuals. On each chart, take
$h_\alpha\in Q_E(U_\alpha)$ with $e_\alpha=h_\alpha+\beta(p_\alpha)$. Descending
\eqref{eq:2.31} to the action of the quotient and computing the difference on the
overlaps gives
\begin{equation*}
r_E=\kappa^1(r_{\mathrm{sem}})+d^0_Eh
\tag{2.37}\label{eq:2.37}
\end{equation*}
Hence the two residuals correspond to the same cohomology class.

\end{proof}

When the two atlases are aligned by $e_\alpha=\beta(p_\alpha)$, the $h$ in
\eqref{eq:2.37} is zero. Even when they are chosen separately, the difference is an
explicit coboundary, so the correspondence of obstruction classes is preserved. This is
the comparison of
\cite[\href{https://arxiv.org/html/2608.21458v1\#S5}{Theorem~5.1(i)--(ii),
\S\S5.3--5.5}]{SAGA}.

\begin{corollary}[SAGA and global repair]\label{cor:2.50}

In addition to the conditions of Theorem~\ref{thm:2.49}, suppose that the actual repair
states $P_{\mathrm{sem}}$ form a sheaf on the chosen site. Then
\begin{equation*}
P_{\mathrm{sem}}(W)\ne\varnothing
\quad\Longleftrightarrow\quad
[r_{\mathrm{sem}}]=0
\quad\Longleftrightarrow\quad
[r_E]=0.
\tag{2.38}\label{eq:2.38}
\end{equation*}

\end{corollary}

\begin{proof}

The first equivalence is obtained by applying Theorem~\ref{thm:2.22} to the torsor
action derived from the repair words. The second follows from the isomorphism of
Theorem~\ref{thm:2.49} and the correspondence of residual classes. Concretely, finding a
correction with $r_{\mathrm{sem}}=d^0b$ and gluing $(-b_\alpha)+p_\alpha$ by the sheaf
condition yields a global repair.

\end{proof}

This is the global repair statement of
\cite[\href{https://arxiv.org/html/2608.21458v1\#S5.SS6}{Theorems~5.1(iii) and~5.2,
\S5.6}]{SAGA}. The zero class expressing the existence of corrections and the gluing of
the repair states themselves are tied together by the sheaf condition.

\begin{example}[Two presentations by parity]\label{ex:2.51}

We use the four-chart cover of Lemma~\ref{lem:2.38}. On the semantic side, we declare on
each nonempty intersection a repair generator $\sigma$ and the relation $2\sigma=0$. On
the equation side, we take the ring $\mathbb{Z}$, the symbolic generator 2, and one
required residual equal to 1. The coefficients on the two sides are obtained as
\begin{equation*}
M_{\mathrm{sem}}=\mathbb{Z}\sigma/(2\sigma),
\qquad Q_E=\mathbb{Z}/(2),
\qquad\chi(\sigma)=[1]
\tag{2.39}\label{eq:2.39}
\end{equation*}
The restrictions between intersections are identity maps. The coefficients are constant
presheaves on nonempty contexts and zero on the empty context. The action of the
semantic repair words is determined by parity alone, and its stabilizer subgroup is
$2\mathbb{Z}\sigma$. The kernel of the map $n\sigma\mapsto[n]$ is $2\mathbb{Z}\sigma$
and its image is all of $\mathbb{Z}/(2)$, so the two completeness conditions can be
checked directly.

The local states can also be constructed. On the semantic side we place the two states
$\{0,1\}$ with changes by parity, and on the equation side the two residue classes
$\{[0],[1]\}$ with the additive action. Orient the four edges as
$a_0\to a_1\to b\to c\to a_0$, and on both sides choose the comparison that shifts the
reference by one on the first edge only and leaves it unshifted on the other three.
Since the coefficients have characteristic 2, returning the last edge to the
orientation of the total order does not change the values. As in
Proposition~\ref{prop:2.41}, these define a sheaf of states obtained by gluing chart
states along affine transitions. The map sending a semantic state $n$ to $[n]$ on the
equation side is $\beta$, and it commutes with restriction and with the action of
generators.

Choosing the zero state on every chart, the residuals on the two sides become
$(1,0,0,0)$ in their respective presentations. The sum over the four edges is 1, while
the sum of any coboundary is zero, so both obstruction classes are nonzero.
Theorem~\ref{thm:2.49} matches these two nonzero classes, and Corollary~\ref{cor:2.50}
shows that no global repair exists under the fixed transitions. If the first transition
is also changed to zero, the zero states glue as they are, and the obstruction classes
on both sides become zero.

\end{example}

\section*{Summary of the Chapter}

In this chapter we connected equations, comparison of local states, and semantic repair
as follows.

\begin{itemize}
\item \textbf{Geometry of Laws.} We built local algebras from the structural relations
and ideals from the symbolic equations. Under the algebraic presentation of the
evaluation data and the localization condition, simultaneous vanishing of the residuals
coincides with factorization through the lawful locus (Theorem~\ref{thm:2.9}).
\item \textbf{Obstructions to local consistency.} The differences on overlaps form a
cocycle, and changing the local states adds a coboundary. When the corrections are
realized as a torsor action and the states form a sheaf, the zero class is equivalent to
the existence of a global state (Theorem~\ref{thm:2.22}).
\item \textbf{Roles of covers and coefficients.} On a forest, obstructions vanish by
surjectivity of the restrictions, while a nonzero period along a cycle detects a
concrete obstruction. In the comparison that removes the structural phase, the
differential preserving the structural part and the vanishing of first-degree
obstructions on the structural side give injectivity
(Propositions~\ref{prop:2.26}, \ref{prop:2.28}, and~\ref{prop:2.32}).
\item \textbf{Concrete local data.} From an Atom family with points and generators we
built covers and an integer-coefficient sheaf, and obtained $o_q(x)=\xi+d^0p$ as the
comparison of lawful chart states along affine transitions
(Proposition~\ref{prop:2.41}).
\item \textbf{Correspondence with the semantics of repair.} The two coefficients, one
built from a presentation by repair words and the other from the quotient by the
equations, become isomorphic under the correspondence of local states and the
completeness conditions. This isomorphism matches the specified obstruction classes and,
combined with the sheaf condition on the actual repair states, characterizes the
existence of a global repair (Theorem~\ref{thm:2.49}, Corollary~\ref{cor:2.50}).
\end{itemize}

For the three services at the opening of the chapter, the correction summed around the
cycle expresses the discrepancy that remains however the reference of each part is
rechosen. If the obstruction class is zero, the local corrections can be applied so that
the states agree on the overlaps. That the corrected states form a single global state
is guaranteed by the sheaf condition on the actual states.

In the next chapter we study what is preserved in the diagnosis when the resolution or
the cover reading the same Laws is changed. We map the concrete obstruction classes
obtained in \S\ref{sec:2.8} into the diagnosis generated from Law evaluations, and
clarify the conditions under which the zero and nonzero information is preserved.
